\subsection{Определение линейного пространства. Простейшие свойства. Примеры}

\begin{definition}
Множество $V$ называется \textbf{линейным (векторным) пространством} над числовым полем $F$ (где $F$ — это $\R$ или $\mathbb{C}$), если на нём введены две операции:
\begin{enumerate}
    \item \textbf{Сложение} (бинарная операция): каждой паре элементов $a, b \in V$ ставится в соответствие единственный элемент $c \in V$, называемый их суммой ($c = a + b$).
    \item \textbf{Умножение на скаляр} (небинарная операция): каждому элементу $a \in V$ и каждому скаляру $\lambda \in F$ ставится в соответствие единственный элемент $c \in V$, называемый произведением вектора на число ($c = \lambda a$).
\end{enumerate}
При этом указанные операции должны удовлетворять восьми аксиомам:
\begin{enumerate}
    \item $a + b = b + a$ для любых $a, b \in V$ (коммутативность сложения).
    \item $(a + b) + c = a + (b + c)$ для любых $a, b, c \in V$ (ассоциативность сложения).
    \item Существует такой элемент $\mathbf{0} \in V$, что $a + \mathbf{0} = a$ для любого $a \in V$ (существование нулевого вектора).
    \item Для любого $a \in V$ существует такой элемент $-a \in V$, что $a + (-a) = \mathbf{0}$ (существование противоположного элемента).
    \item $1 \cdot a = a$ для любого $a \in V$, где $1$ — единица поля $F$.
    \item $\lambda(\mu a) = (\lambda\mu)a$ для любых $\lambda, \mu \in F$ и $a \in V$.
    \item $(\lambda + \mu)a = \lambda a + \mu a$ для любых $\lambda, \mu \in F$ и $a \in V$.
    \item $\lambda(a + b) = \lambda a + \lambda b$ для любого $\lambda \in F$ и $a, b \in V$.
\end{enumerate}
\end{definition}

В лекциях элементы пространства $V$ называются \textbf{векторами} (обозначаются латинскими буквами без стрелок), а элементы поля $F$ — \textbf{скалярами} (обозначаются греческими буквами).

\begin{theorem}[Простейшие свойства]
В любом линейном пространстве выполняются следующие свойства:
\begin{enumerate}
    \item Нулевой вектор $\mathbf{0}$ единственен.
    \item Для каждого вектора $a$ противоположный вектор $-a$ единственен.
    \item $0 \cdot a = \mathbf{0}$ для любого $a \in V$ (здесь $0$ — числовой ноль, $\mathbf{0}$ — векторный).
    \item $(-1) \cdot a = -a$ для любого $a \in V$.
    \item $\lambda \cdot \mathbf{0} = \mathbf{0}$ для любого скаляра $\lambda \in F$.
\end{enumerate}
\end{theorem}

\begin{tcolorbox}[title=Доказательство, breakable]
\textbf{1. Единственность нуля.}
Предположим, что существуют два нуля $\mathbf{0}_1$ и $\mathbf{0}_2$. По аксиоме (3):
\[ \mathbf{0}_1 = \mathbf{0}_1 + \mathbf{0}_2 \]
С другой стороны, в силу коммутативности и той же аксиомы:
\[ \mathbf{0}_2 = \mathbf{0}_2 + \mathbf{0}_1 = \mathbf{0}_1 + \mathbf{0}_2 \]
Следовательно, $\mathbf{0}_1 = \mathbf{0}_2$.

\textbf{2. Единственность противоположного элемента.}
Пусть для вектора $a$ существуют два противоположных элемента $a'_1$ и $a'_2$. Тогда:
\[ a'_1 = a'_1 + \mathbf{0} = a'_1 + (a + a'_2) = (a'_1 + a) + a'_2 = \mathbf{0} + a'_2 = a'_2 \]

\textbf{3. Умножение числового нуля на вектор.}
Рассмотрим равенство $0 + 0 = 0$ в поле $F$. Тогда $(0 + 0)a = 0a$. По аксиоме (7): $0a + 0a = 0a$. Прибавим к обеим частям противоположный элемент для $0a$ (обозначим его $-(0a)$):
\[ 0a + 0a + (-(0a)) = 0a + (-(0a)) \implies 0a + \mathbf{0} = \mathbf{0} \implies 0a = \mathbf{0} \]

\textbf{4. Связь противоположного вектора с числом $-1$.}
Покажем, что сумма $a + (-1)a$ дает $\mathbf{0}$:
\[ a + (-1)a = 1 \cdot a + (-1) \cdot a = (1 + (-1))a = 0 \cdot a = \mathbf{0} \]
В силу единственности противоположного элемента, $(-1)a = -a$.

\textbf{5. Умножение скаляра на нулевой вектор.}
Рассмотрим $\mathbf{0} + \mathbf{0} = \mathbf{0}$. Тогда $\lambda(\mathbf{0} + \mathbf{0}) = \lambda\mathbf{0}$. По аксиоме (8): $\lambda\mathbf{0} + \lambda\mathbf{0} = \lambda\mathbf{0}$. Прибавляя противоположный к $\lambda\mathbf{0}$, получаем $\lambda\mathbf{0} = \mathbf{0}$.
\end{tcolorbox}

\textbf{Примеры линейных пространств:}
\begin{enumerate}
    \item \textbf{Геометрические векторы ($LPV$):} Множество векторов на прямой, на плоскости или в трехмерном пространстве с обычными операциями сложения и умножения на число.
    \item \textbf{Координатное пространство ($F^n$):} Множество упорядоченных наборов из $n$ чисел (строк или столбцов). Операции проводятся покомпонентно.
    \item \textbf{Пространство матриц $M_{m,n}$:} Множество прямоугольных матриц размера $m \times n$. Нулем является нулевая матрица.
    \item \textbf{Пространство непрерывных функций $C[a, b]$:} Функции, непрерывные на отрезке $[a, b]$. Это пример бесконечномерного пространства.
    \item \textbf{Пространство многочленов $P_n$:} Многочлены степени не выше $n$. \textit{Важное замечание:} множество многочленов степени \textbf{ровно} $n$ ($n > 0$) не является линейным пространством, так как операция сложения может вывести из этого множества (например, $t^n + (-t^n) = 0$), и в нём нет нулевого элемента.
    \item \textbf{Пространство решений однородной СЛАУ:} Если $X_1$ и $X_2$ — решения системы $AX = \mathbf{0}$, то их сумма и произведение на число также являются решениями. Множество решений неоднородной системы линейным пространством \textbf{не является}.
\end{enumerate}

\begin{tcolorbox}[title=Источники, colback=gray!10]
\begin{itemize}
  \item Лекция №\,4, тема «Линейные пространства», временной отрезок 140:10--154:00
  \item PDF «Линейная алгебра\_compressed (1).pdf», раздел 1 «Понятие линейного пространства», стр. 703--707
  \item PDF «Вопросы», вопрос №\,10
\end{itemize}
\end{tcolorbox}
